\chapter{Complexity — Input Length}
%%% generated by the blueprint pipeline from TCSlib.Complexity.TuringMachine.Build.InputLength
\section{Overview}

Declarations of the Lean module \texttt{TCSlib.Complexity.TuringMachine.Build.InputLength}.

\section{Declarations}

\begin{definition}[Input-length machine phases]
\lean{Turing.InputLengthPhase}
\label{Turing.InputLengthPhase}
\leanok
The input-length machine has six control phases: counting the input, propagating a
carry, moving back, beginning the return scan, continuing the return scan, and
completion.
\end{definition}

\begin{definition}[Input-length counting machine]
\lean{Turing.inputLengthMachine}
\label{Turing.inputLengthMachine}
\leanok
\uses{Turing.Action.mapState, Turing.FinTM.controlAction, Turing.InputLengthPhase, Turing.MultiTapeTM, Turing.inputFoldTM, Turing.inputMoveAction, Turing.tapeAction}
For an input over an arbitrary symbol alphabet, the input-length machine is a
deterministic machine with one work tape over the Boolean alphabet and the six phases
for counting, carry propagation, moving back, returning to the input, scanning back
across it, and completion. Its transition function performs the corresponding input and
work-tape actions in each phase.
\end{definition}

\begin{theorem}[Binary Digits Under Increment]
\lean{Turing.inputLength_digits}
\label{Turing.inputLength_digits}
\leanok
\difficulty{5}
For every natural number $n$, there are a natural number $p$ and a Boolean list $v$ such
that either the binary digits of $n$ consist of $p$ ones followed by a zero and then
$v$, or they consist of $p$ ones with $v$ empty; the binary digits of $n+1$ then consist
of $p$ zeros followed by a one and then $v$.
\end{theorem}

\begin{theorem}[Binary Digit Charge]
\lean{Turing.inputLength_charge}
\label{Turing.inputLength_charge}
\leanok
\difficulty{3}
Let $n,p$ be natural numbers and let $v$ be a Boolean list. Suppose either the binary
digits of $n$ consist of $p$ ones followed by a zero and then $v$, or they consist of
$p$ ones with $v$ empty; and suppose the binary digits of $n+1$ consist of $p$ zeros
followed by a one and then $v$. Then the number of ones in the binary digits of $n+1$,
plus $p$, equals the number of ones in the binary digits of $n$ plus one, and $p$ is
strictly less than the number of binary digits of $n+1$.
\end{theorem}

\begin{definition}[Input-length configuration update]
\lean{Turing.inputLengthCfg}
\label{Turing.inputLengthCfg}
\leanok
\uses{Turing.Cfg, Turing.FinTM.bufferTape, Turing.InputLengthPhase}
Let $x$ and $w$ be finite Boolean strings, let $d$ be a configuration of the
one-work-tape input-length machine on input $x$, let $r$ be a natural number, let $s$ be
one of its phases, let $ip$ be a valid input-head position for $x$, and let $z$ be an
integer. The updated configuration has state $s$, input-head position $ip$, and output
equal to that of $d$; its work-head position is shifted by $z$, and on the $r$ cells
beginning at the old work-head position its tape contains $w$, while all other cells
retain their contents from $d$.
\end{definition}

\begin{theorem}[Zero-Length Tape Overlay]
\lean{Turing.inputLengthCfg_zero}
\label{Turing.inputLengthCfg_zero}
\leanok
\uses{Turing.Cfg, Turing.InputLengthPhase, Turing.inputLengthCfg}
\difficulty{4}
Let $x$ and $w$ be Boolean words, let $d$ be a one-work-tape configuration on input $x$,
and choose a phase $s$ and input-head position $ip$. The configuration obtained by
overlaying a word of length zero on the work tape, with no work-head displacement, is
exactly $d$ with only its state set to $s$ and its input-head position set to $ip$.
\end{theorem}

\begin{theorem}[Writing a Buffered Tape Symbol]
\lean{Turing.inputLength_write}
\label{Turing.inputLength_write}
\leanok
\uses{Turing.Cfg, Turing.Cfg.inputSymbol, Turing.Cfg.workTapeSymbols, Turing.InputLengthPhase, Turing.MultiTapeTM.step, Turing.catalog_write_take, Turing.inputLengthCfg, Turing.inputLengthMachine, Turing.moveInputPos_zero, Turing.step_of_tapeAction_write, Turing.tapeAction}
\difficulty{4}
Let $x,w$ be Boolean words, let $d$ be a one-work-tape configuration on input $x$, let
$r$ be a natural number with $r<|w|$, and choose phases $s,s'$, an input-head position
$ip$, and a sign $m\in\{-1,0,1\}$. If the machine’s transition at phase $s$, when given
the input and work-tape symbols in the configuration with $w$ buffered through its first
$r$ symbols and the work head at offset $r$, is the action that writes the $r$-th symbol
of $w$ on that tape, moves its head by $m$, leaves the input head still, produces no
output, and enters phase $s'$, then one machine step yields the configuration with the
first $r+1$ symbols of $w$ buffered, phase $s'$, unchanged input-head position, and
work-head displacement $r+m$ from its original position.
\end{theorem}

\begin{theorem}[One Input-Length Counting Round]
\lean{Turing.inputLength_round}
\label{Turing.inputLength_round}
\leanok
\uses{Turing.Cfg, Turing.Cfg.inputSymbol, Turing.Cfg.workTapeSymbols, Turing.FinTM.bufferTape, Turing.FinTM.bufferTape_left, Turing.FinTM.bufferTape_nat, Turing.InputLengthPhase, Turing.MultiTapeTM.runFrom, Turing.MultiTapeTM.runFrom_succ_eq_step', Turing.MultiTapeTM.step, Turing.inputLengthCfg, Turing.inputLengthCfg_zero, Turing.inputLengthMachine, Turing.inputLength_charge, Turing.inputLength_write, Turing.inputSymbolInner, Turing.moveInputPos, Turing.moveInputPos_zero, Turing.step_of_inputMoveAction, Turing.step_of_tapeAction_move}
\difficulty{6}
Let $x$ be a Boolean word and let $d$ be a one-work-tape configuration on input $x$,
with natural numbers $n,p$ and a Boolean list $v$. Suppose $d$ is in phase count, its
input head is at position $n+1$, $n<|x|$, the binary digits of $n$ have either the form
of $p$ ones followed by a zero and then $v$, or $p$ ones with $v$ empty, and the binary
digits of $n+1$ are $p$ zeros followed by a one and then $v$. Suppose also that the work
tape agrees with the buffer tape for the binary digits of $n$ at every offset from the
work-head position between $-1$ and the number of binary digits of $n$. After $2p+3$
steps, the machine is in phase count with the binary digits of $n+1$ stored over a
segment of length $p+1$, the input head moved by one clamped step to the right, and the
work head left at its original position. At each time $u\le 2p+3$, the work-head
position is between one cell left of its initial position and its initial position plus
$p$.
\end{theorem}

\begin{theorem}[Reading the Buffered Input]
\lean{Turing.inputLengthCfg_read}
\label{Turing.inputLengthCfg_read}
\leanok
\uses{Turing.Cfg, Turing.FinTM.bufferTape, Turing.InputLengthPhase, Turing.inputLengthCfg}
\difficulty{4}
Let $x$ and $w$ be Boolean words, let $d$ be a one-work-tape configuration on input $x$,
and let $s$ be a machine phase and $ip$ an input-head position. Suppose the original
work tape is blank at every cell $d$’s work-head position plus $b$, for all integers $b$
with $-1\le b\le N$. Then, for every integer $a$ with $-1\le a\le N$, the configuration
formed by overlaying the entire word $w$ on that tape and leaving the work head in place
reads at cell $d$’s work-head position plus $a$ exactly the symbol that the buffer tape
for $w$ stores at offset $a$.
\end{theorem}

\begin{theorem}[Merging configurations for consecutive counts]
\lean{Turing.inputLengthCfg_merge}
\label{Turing.inputLengthCfg_merge}
\leanok
\uses{Turing.Cfg, Turing.FinTM.bufferTape, Turing.InputLengthPhase, Turing.inputLengthCfg, Turing.inputLength_charge}
\difficulty{6}
Let $x$ be a Boolean word and $d$ a configuration on input $x$. Let $n,p$ be natural
numbers and $v$ a Boolean word. Suppose either the binary expansion of $n$ consists of
$p$ ones followed by a zero and then $v$, or it consists of exactly $p$ ones and $v$ is
empty; also suppose the binary expansion of $n+1$ consists of $p$ zeros followed by a
one and then $v$. For any two valid input-head positions $ip,ip'$, applying the
configuration construction first with the binary expansion of $n$ and then with that of
$n+1$ gives the same configuration as applying it directly to the binary expansion of
$n+1$.
\end{theorem}

\begin{theorem}[Input-Length Counting Run]
\lean{Turing.inputLength_count}
\label{Turing.inputLength_count}
\leanok
\uses{Turing.Cfg, Turing.FinTM.bufferTape, Turing.InputLengthPhase, Turing.MultiTapeTM.runFrom, Turing.MultiTapeTM.runFrom_add, Turing.MultiTapeTM.runFrom_zero, Turing.inputLengthCfg, Turing.inputLengthCfg_merge, Turing.inputLengthCfg_read, Turing.inputLengthCfg_zero, Turing.inputLengthMachine, Turing.inputLength_charge, Turing.inputLength_digits, Turing.inputLength_round, Turing.moveInputPos, Turing.moveInputPos_pos_of_ne_right}
\difficulty{6}
Let $x$ be a Boolean word and let $d$ be a one-work-tape configuration on input $x$, in
phase count and with its input head at position $1$. Suppose the work tape is blank at
offsets from the work-head position in the range $-1$ through the number of binary
digits of the length of $x$. For every natural number $n\le |x|$, there is a time $t$
such that $t+2c(n)=5n$, where $c(n)$ is the number of ones in the binary expansion of
$n$, and after $t$ steps the machine is in phase count with the input head at position
$n+1$, the work tape containing the binary expansion of $n$ over a segment of its
length, and the work head at its original position. At every time $u\le t$, the
work-head position lies between one cell left of its initial position and its initial
position plus the number of binary digits of $n$.
\end{theorem}

\begin{theorem}[Machine computing input length]
\lean{Turing.exists_inputLengthTM}
\label{Turing.exists_inputLengthTM}
\leanok
\uses{Turing.Cfg, Turing.Cfg.inputSymbol, Turing.FinTM.bufferTape, Turing.FinTM.rewind_heads, Turing.InputLengthPhase, Turing.MultiTapeTM, Turing.MultiTapeTM.runFrom, Turing.MultiTapeTM.runFrom_add, Turing.MultiTapeTM.runFrom_succ_eq_step', Turing.MultiTapeTM.step, Turing.inputLengthCfg, Turing.inputLengthMachine, Turing.inputLength_count, Turing.moveInputPos_zero, Turing.step_of_inputMoveAction}
\difficulty{7}
There exist a natural number $K$, a finite state set, a deterministic Turing machine
over the binary alphabet with $K+1$ work tapes, entry and exit states, and a constant
$C$. From any configuration on a binary input $x$ with the machine in the entry state,
input head at position $1$, tape $0$ blank on the interval from one cell before its head
through the binary representation of $|x|$, and every other work tape entirely blank,
the machine reaches its exit state at some time $t \le C(|x|+1)$. It does not reach the
exit state earlier, every work-tape head stays within distance $C(|x|+1)$ of its initial
position through time $t$, and at time $t$ tape $0$ holds the binary representation of
$|x|$ starting at its initial head position, with all other cells and configuration data
unchanged.
\end{theorem}
